%%%%%%%%%%%%%%%%%%%%%%%%%%%%%%%%%%%%%%%%%%%%%%%%%%%%%%%%%%%%%%%%%%%%%%%%%%%
% Handout und Uebungsblatt zu Lehreinheit 1
% Im Skript am Kapitelende, in handouts.tex als Druckvorlage.
%%%%%%%%%%%%%%%%%%%%%%%%%%%%%%%%%%%%%%%%%%%%%%%%%%%%%%%%%%%%%%%%%%%%%%%%%%%
\begin{handout}{Staunen als Erkenntnishaltung}

\textbf{Worum es heute ging.} Ein Pendel schwingt über Ihrer Hand, obwohl
Sie \enquote{nichts tun}. Mit geschlossenen Augen wird der Ausschlag
kleiner, mit aufgelegtem Unterarm verschwindet er fast. Sie selbst bewegen
das Pendel -- mit Muskelimpulsen unterhalb Ihrer Wahrnehmungsschwelle
(ideomotorischer Effekt).

\textbf{Kernbegriffe}
\begin{description}
  \item[Prediction Error] Abweichung zwischen dem, was das Gehirn erwartet,
    und dem, was die Sinne melden. Große Abweichungen erleben wir als
    Überraschung; sie sind der Anlass, ein Modell zu ändern, also zu lernen.
  \item[Predictive Processing] Modell der Wahrnehmung, nach dem das Gehirn
    fortlaufend Vorhersagen erzeugt und vor allem die Abweichungen
    verarbeitet \parencite{clark2016}.
  \item[Belief] bewusste Überzeugung, durch Argumente änderbar.
  \item[Alief] automatische, körpernahe Reaktionsbereitschaft, die neben
    dem Belief weiterläuft \parencite{gendler2008}.
  \item[Ideomotorischer Effekt] die Vorstellung einer Bewegung erzeugt
    minimale, unbemerkte Muskelbewegungen \parencite{carpenter1852}.
\end{description}

\textbf{Der Dreiklang des Semesters}

\begin{center}
\begin{tabular}{p{0.28\textwidth}p{0.28\textwidth}p{0.32\textwidth}}
\toprule
\textbf{Staunen} & \textbf{Erklärung} & \textbf{Methode} \\
\midrule
Das Modell passt nicht. & Hypothesen sammeln, noch ohne Bewertung. & Ein Verfahren finden, das zwischen den Hypothesen entscheidet. \\
\bottomrule
\end{tabular}
\end{center}

\textbf{Drei historische Kontrollen}
\begin{itemize}
  \item \emph{Chevreul (1833):} Glasplatte, Armstütze, Augenbinde -- die
    Ursache liegt in der Hand.
  \item \emph{Carpenter (1852):} benennt das ideomotorische Prinzip.
  \item \emph{Faraday (1853):} ein Zeiger macht das unbewusste Schieben beim
    Tischrücken sichtbar; wer den Zeiger sieht, schiebt nicht mehr.
\end{itemize}

\begin{merkbox}
Das Gefühl, recht zu haben, ist kein Befund. Wer dem eigenen Erleben
misstraut, leugnet es nicht, sondern misst es.
\end{merkbox}

\textbf{Ihr Semesterritual: das Tagebuch des Staunens.} Notieren Sie jede
Woche ein Erlebnis, das Sie überrascht hat. Halten Sie fest: (1) Was habe
ich erwartet? (2) Was ist passiert? (3) Welcher Begriff aus der Vorlesung
passt dazu? Das Tagebuch ist Teil der Portfolioprüfung (20\,\%).

\textbf{Zum Weiterlesen:} \textcite{seth2021},
Kap.~1--4; \textcite{keltner2003}.
\end{handout}

\begin{uebungsblatt}{Staunen, Erklärung, Methode}

\aufgabe{Hypothesen sortieren}[in der Stunde, Kleingruppe, 20 Min.]
Übertragen Sie die Erklärungen von der Tafel in die Tabelle. Ordnen Sie
jede einer Stufe zu und geben Sie für jede Erklärung eine Beobachtung an,
die sie \emph{widerlegen} würde.

\begin{center}
\renewcommand{\arraystretch}{2.2}
\begin{tabular}{|p{0.38\textwidth}|p{0.14\textwidth}|p{0.38\textwidth}|}
\hline
\textbf{Aussage / Erklärung} & \textbf{Stufe} & \textbf{Was würde sie widerlegen?} \\
\hline
 & & \\ \hline
 & & \\ \hline
 & & \\ \hline
 & & \\ \hline
\end{tabular}
\end{center}

\aufgabe{Die beiden Kontrollen deuten}[in der Stunde, Einzelarbeit]
\begin{enumerate}
  \item Was haben Sie beobachtet, als Sie die Augen geschlossen haben?
    Was folgt daraus für die Ursache der Bewegung?
  \item Was haben Sie beobachtet, als der Unterarm auf dem Tisch lag?
\end{enumerate}
\schreiblinien{4}

\aufgabe{Belief oder Alief?}[Hausaufgabe]
Nennen Sie zwei eigene Beispiele, in denen Ihr Wissen und Ihr Erleben
auseinanderfallen (Beispiel: Sie wissen, dass eine Glasbrücke trägt, und
trauen sich trotzdem nicht darauf). Welches Verfahren könnte Ihnen helfen,
dem Wissen statt dem Erleben zu folgen?
\schreiblinien{4}

\aufgabe{Kurzessay}[Hausaufgabe, 500 Wörter, Abgabe bis LE~2]
\enquote{Beschreiben Sie ein Erlebnis, das Sie zum ersten Mal an einer
Überzeugung zweifeln ließ. Was hat das mit Wissenschaft zu tun?} Verwenden
Sie mindestens zwei der Kernbegriffe dieser Einheit.

\aufgabe{Vorbereitung auf LE~2}[Hausaufgabe]
Bringen Sie eine Frage mit, die Sie \emph{wirklich} beschäftigt -- aus
Studium, Beruf oder Alltag. Sie wird in LE~2 zur Forschungsfrage geschärft.

\begin{lehrhinweis}[Erwartungshorizont]
Aufgabe~1: Mindestens eine Erklärung je Stufe; Widerlegungsbedingungen
müssen beobachtbar sein (\enquote{das Pendel schwingt auch bei
aufgelegtem Arm} statt \enquote{es stimmt nicht}). Aufgabe~2: Augen zu
$\to$ visuelle Rückkopplung fehlt, Ausschlag sinkt; Arm auf $\to$
Muskelimpulse werden abgefangen. Beides spricht für die Hand als Ursache.
Aufgabe~3: Beispiele mit Höhenangst, Placebo, optischen Täuschungen;
Verfahren = Messung, Protokoll, Blindbedingung.
\end{lehrhinweis}
\end{uebungsblatt}
